\renewcommand{\Anno}{2004}

\renewcommand{\NumeroProblema}{2}

\renewcommand{\Testo}{È ben noto che in Svizzera la benzina costa
meno che in Italia, e che il suo prezzo sale man mano che ci si avvicina alla frontiera.  Un automobilista, diretto in Italia, fa il pieno a Zurigo e decide che nel tragitto in Svizzera si fermerà al più una volta per fare rifornimento.  Arrivato alla frontiera italiana farà poi nuovamente il pieno, al prezzo italiano.
	
Supponendo che il consumo di benzina dell'automobile sia costante, che la capacità del serbatoio sia tale da consentire di completare il percorso anche senza soste, che ci sia un distributore in ogni punto del percorso, e che il prezzo al litro aumenti linearmente avvicinandosi alla frontiera italiana, determinare il punto
economicamente più favorevole in cui fermarsi.}

\renewcommand{\Risposta}{Conviene fare una sosta a metà del percorso.}

\renewcommand{\Aiutino}{Occorre matematizzare il problema.
Consideriamo il percorso come un segmento di estremi $Z$ ed $I$ e lunghezza $l$.  Indichiamo con $c$ il consumo per unità di lunghezza, con $P_{Z}$ e $P_{I}$ i prezzi a Zurigo ed in Italia, e con $x$ la distanza da $Z$ del punto in cui si fa la sosta ($x=0$ corrisponde a non rifornire).  Calcolato il prezzo in $x$, si ha che il costo totale dell'operazione è
$$c\left\{(P_{I}-P_{Z})\left(\frac{x^{2}}{l}-x\right)+lP_{I}\right\}.$$

Il minimo di tale espressione si ha per $x=\ldots$}

\renewcommand{\SoluzioneUfficiale}{\TitoloSoluzione Indichiamo con $l$ 
la lunghezza del percorso, con $c$ il consumo in litri per unità di lunghezza, con $P_{Z}$ e $P_{I}$, rispettivamente, i prezzi in euro di un litro di benzina a Zurigo ed al confine italiano. Indichiamo con $x$ la distanza da Zurigo del punto in cui si fa la sosta.

Determiniamo ora il prezzo $P(x)$ di un litro di benzina in funzione di $x$. Per ipotesi sappiamo che deve essere una funzione del tipo $P(x)=mx+n$ che deve valere $P_{Z}$ per $x=0$ e $P_{I}$ per $x=l$. Queste condizioni si traducono nel sistema
$$n=P_{Z},\hspace{4em}ml+n=P_{I},$$
da cui
$$P(x)=\frac{P_{I}-P_{Z}}{l}\,x+P_{Z}.$$

Quando giunge in $x$, l'automobilista ha consumato $cx$ litri di carburante, che rabbocca al prezzo di $P(x)$ euro al litro; nella rimanente parte del viaggio consumerà ancora $c(l-x)$ litri che poi rifornirà in Italia al prezzo di $P_{I}$ euro al litro. Il costo totale dell'operazione risulta pertanto
$$cxP(x)+c(l-x)P_{I}=c\left\{(P_{I}-P_{Z})
\left(\frac{x^{2}}{l}-x\right)+lP_{I}\right\}.$$

Si noti che il caso $x=0$ corrisponde a non fare nessuna sosta lungo il percorso. Poiché il coefficiente di $x^{2}$ è positivo, il grafico di questa funzione della $x$ è una parabola con la concavità rivolta verso l'alto, che assume valori uguali per $x=0$ e $x=l$ ed ha minimo nel vertice, corrispondente ad $x=l/2$.

Ne segue che conviene fare una sosta a metà del percorso,
indipendentemente da tutti gli altri parametri.}

\renewcommand{\SoluzioniAlternative}{
}

\renewcommand{\Osservazioni}{\TitoloOsservazione Ci sono ovviamente 
molti metodi per determinare l'ascissa del vertice di una parabola: utilizzare le apposite formule della geometria analitica, oppure ricorrere al calcolo differenziale (basta vedere dove si annulla la derivata prima), oppure sfruttare la simmetria come abbiamo fatto nella soluzione: se un trinomio di secondo grado assume lo stesso valore in 2 punti $x_{1}$ e $x_{2}$, allora l'ascissa del vertice è il punto medio $(x_{1}+x_{2})/2$.
}

\renewcommand{\Commenti}{
}